% from ega4-9-fr.tex
\oldpage[IV-3]{87}

\begin{corollary}[9.8.7]
\label{IV.9.8.7-fr}
Soient $f:X\to S$ un morphisme de présentation finie, $\mathcal F$ un
$\mathcal O_X$-Module quasi-cohérent de présentation finie. Pour tout $s\in S$,
soit $M(s)$ l'ensemble des couples $(d,m)$ tels qu'il existe une composante
irréductible de $\operatorname{Supp}(\mathcal F_s)$ de dimension $d$ et de
multiplicité radicielle $m$ pour $\mathcal F_s$ (4.7.8). Alors la fonction
$s\mapsto M(s)$ est localement constructible dans $S$.
\end{corollary}

Il résulte de (4.2.7) et (4.7.9) que la condition (9.2.1, (i)) est satisfaite
pour la propriété dont nous voulons démontrer qu'elle est constructible. On est
donc ramené au cas où $S$ est affine, noethérien et intègre et à prouver que $M$
est constante dans un voisinage du point générique de $S$~; mais alors la
proposition résulte de (9.8.3) et (9.8.6).

\begin{proposition}[9.8.8]
\label{IV.9.8.8-fr}
Soient $f:X\to S$ un morphisme de présentation finie, $\mathcal F$ un
$\mathcal O_X$-Module quasi-cohérent de présentation finie. Pour tout $s\in S$,
soit $l(s)$ la somme des multiplicités totales pour $\mathcal F_s$ (relatives à
$\boldsymbol{k}(s)$) des points génériques des composantes irréductibles de
$\operatorname{Supp}(\mathcal F_s)$ (4.7.12). Alors la fonction $s\mapsto l(s)$
est localement constructible dans $S$.
\end{proposition}

Compte tenu de (4.7.12), la condition (9.2.1, (i)) est satisfaite pour la
propriété dont nous voulons démontrer qu'elle est constructible, et on est donc
ramené au cas où $S$ est affine, noethérien et intègre de point générique
$\eta$, et à montrer que $l$ est constante dans un voisinage de $\eta$.
Utilisant les notations de (9.8.2), cela résulte de la définition (4.7.12), du
fait que le nombre géométrique des composantes irréductibles de chaque
$(Z_i)_s$ est constant dans un voisinage de $\eta$ (9.7.8), que la longueur
géométrique de $\mathcal F_s$ en $x_{is\alpha}$ est également à celle de
$\mathcal F_\eta$ en $x_i$ pour chaque $i$ tel que $(Z_i)_\eta$ soit maximal
(9.8.6) et enfin de ce que l'adhérence de $\{x_{is\alpha}\}$ est un cycle
premier associé maximal de $\mathcal F_s$ si et seulement si $(Z_i)_\eta$ est
un cycle premier associé maximal de $\mathcal F_\eta$ (9.8.3).

\begin{remark}[9.8.9]
\label{IV.9.8.9-fr}
On peut préciser de diverses façons les propositions précédentes~;
bornons-nous à un énoncé à titre d'exemple. Disons qu'une partie finie $P$ d'un
$k$-préschéma algébrique $X$ est \emph{saturée} si, pour tout couple de points
$x$, $y$ de $P$, les points génériques des composantes irréductibles de
$\overline{\{x\}}\cap\overline{\{y\}}$ appartiennent aussi à $P$~; pour toute
partie finie $Q$ de $X$, il existe une plus petite partie finie $P$ de $X$
contenant $Q$ et saturée~; nous dirons que $P$ est le \emph{saturé} de $Q$. Pour
tout $\mathcal O_X$-Module cohérent $\mathcal F$, nous appellerons
\emph{squelette primaire} de $\mathcal F$ le système $(P,Q,\omega,d,m)$ où
$Q=\operatorname{Ass}(\mathcal F)$, $P$ est le saturé de $Q$, $\omega$ la
relation d'ordre $\overline{\{x\}}\subset\overline{\{y\}}$ sur $P$, $d$ la
fonction $x\mapsto\dim\overline{\{x\}}$ sur $P$, $m$ la fonction
$x\mapsto\operatorname{long}_{\mathcal O_x}\mathcal F_x$ définie sur l'ensemble
des éléments de $Q$ maximaux pour la relation $\omega$. Appelons d'autre part
\emph{squelette virtuel} tout système $(P,Q,\omega,d,m)$, où $P$ est un
ensemble, $Q$ une partie de $P$, $\omega$ une relation d'ordre sur $P$, $d$ une
application de $P$ dans $\mathbf N$, $m$ une application dans $\mathbf N$ de
l'ensemble des éléments maximaux de $Q$~; on définit de façon évidente la notion
d'\emph{isomorphisme} de deux squelettes virtuels. Enfin, avec les notations
précédentes, appelons \emph{type primaire} de $\mathcal F$ la classe (pour la
relation d'isomorphie des squelettes virtuels) du squelette primaire de
$\mathcal F$. Il résulte de (4.2.6), (4.2.7), (4.2.8), (4.5.1) et (4.7.9) que
si, pour une extension algébriquement close $k'$ de $k$, on pose
$X'=X\otimes_k k'$ et $\mathcal F'=\mathcal F\otimes_k k'$, le
\emph{type primaire} de $\mathcal F'$ est indépendant de

